El principal limitador técnico para el despliegue de estos sistemas en servidores propios es el consumo de la memoria de video o VRAM de la tarjeta gráfica. Cada parámetro del modelo debe cargarse en este espacio para que la generación de texto sea rápida y eficiente. Por defecto los modelos se entrenan con una precisión de 16 bits lo cual implica que una arquitectura de pocos parámetros puede requerir una cantidad de memoria inasumible para hardware de rango medio. 

Para superar las limitaciones de hardware, se han implementado diversas estrategias de optimización que actúan sobre la arquitectura y la representación de los datos. La optimización técnica comienza con los formatos de precisión que determinan cuántos bits utiliza el hardware para representar cada decimal de la arquitectura. Al pasar de estándares de dieciséis bits a otros más reducidos como el FP8 se ahorra espacio de forma nativa sin perder apenas precisión \cite{micikevicius2017mixed}. Este método se refuerza con la cuantización de parámetros la cual actúa como una compresión matemática mucho más agresiva mediante la discretización de los pesos y su mapeo sistemático hacia un dominio numérico finito \cite{yadav2025optimizing}. Mientras los formatos de precisión reducen el tamaño estructural del contenedor, la cuantización redefine el contenido para garantizar que un modelo de gran escala quepa en la VRAM de una tarjeta gráfica doméstica.

Para gestionar las limitaciones físicas del hardware se emplea el \textit{offloading} como una técnica logística que distribuye las capas del modelo entre la memoria de vídeo y la memoria RAM del sistema, permitiendo cargar un modelo que demanda más VRAM de la que dispone la tarjeta gráfica. Toda esta infraestructura se beneficia de la arquitectura \textit{Mixture of Experts} o MoE donde el sistema solo activa una pequeña fracción de sus parámetros para responder a cada pregunta concreta \cite{shazeer2017outrageously}. Por ejemplo, los modelos Gemma 4 e4B, Gemma 4 e2B y Qwen 3.6 35B de la tabla \ref{TB:LLM_COMPARATIVA} utilizan esta arquitectura para agilizar la latencia de respuesta preservando la base de conocimiento.